\documentclass[11pt]{article}
\usepackage[a4paper,margin=1in]{geometry}
\usepackage{times}
\usepackage[T1]{fontenc}
\usepackage{amsmath}
\usepackage{graphicx}
\usepackage{booktabs}
\usepackage{array}
\usepackage[colorlinks=true,linkcolor=blue,citecolor=blue,urlcolor=blue]{hyperref}
\usepackage{caption}
\usepackage{setspace}
\usepackage[round,numbers,sort&compress]{natbib}
\bibliographystyle{unsrtnat}

\setstretch{1.08}

\title{\Large Redocking Validation Reveals Structure-Dependent Docking Reliability in ATR: A Reproducible Pilot Virtual Screen of BRICS-Generated Candidates}

\author{[Author name(s) to be added]\\[2pt]
\small [Affiliation to be added]\\[2pt]
\small Correspondence: [email to be added]}
\date{}

\begin{document}
\maketitle

\begin{abstract}
\noindent
\textbf{Background:} Nearly all existing chemical probes of the ATR--CHK1
replication-stress checkpoint are ATP-competitive kinase inhibitors. Two
near-atomic-resolution cryo-EM structures of the human ATR--ATRIP complex
released in 2025 (PDB 9L40, bound to berzosertib/VE-822; PDB 9L4B, bound to
camonsertib/RP-3500) provide new templates for structure-based work on this
target, and 9L40 additionally revealed a second, allosteric berzosertib
binding site at the ATR--ATR dimer interface distinct from the canonical
ATP pocket. Whether newly deposited structures of this kind are equally
reliable for conventional ATP-pocket docking is not something resolution
alone can guarantee. \textbf{Methods:} We assembled a small seed library of
five ATR/CHK1-directed compounds, three extracted directly from the
co-crystallized active-site ligand of PDB 9L40, 9L4B, and a CHK1 reference
structure (PDB 2YM8), generated derivatives by BRICS recombination,
filtered for developability by standard descriptor rules, and --
before docking any candidate -- validated each receptor structure by
redocking its own native ligand against a pre-specified 2.0~{\AA} RMSD pass
threshold, repeated across three independent seeded AutoDock Vina runs.
\textbf{Results:} Redocking succeeded for 9L4B (mean RMSD
$0.61\pm0.02$~{\AA}) and 2YM8 ($0.90\pm0.01$~{\AA}) but reproducibly
\emph{failed} for 9L40 ($2.96\pm0.06$~{\AA} across all three seeds, despite
tight search reproducibility of $0.56$~{\AA} pairwise), so 9L40 was excluded
from candidate scoring per the pre-specified rule. A 40-compound pilot set
(5 seeds + 35 developability-passed BRICS derivatives) was then docked
against the validated 9L4B structure: predicted affinities ranged from
$-9.80$ to $-7.82$~kcal/mol (best pose across three seeded replicates), with
camonsertib itself -- the receptor's own co-crystallized ligand -- ranking
first as expected; no derivative exceeded the best seed compound, and the
top-ranked derivatives were close BRICS analogs of camonsertib sharing its
oxabicyclic-ether/pyrazolopyrazine core with a single substituent varied.
\textbf{Conclusions:} A newly deposited, high-resolution structure is not
automatically a reliable docking receptor, and this can only be established
empirically, per structure, before any candidate is scored -- not assumed
from resolution or recency. All code, data, and structures used are openly
available for independent verification.
\end{abstract}

\noindent\textbf{Keywords:} ATR kinase; molecular docking; redocking
validation; BRICS; virtual screening; reproducibility; AutoDock Vina

\section*{Background}

The ATR--CHK1 checkpoint kinase axis is the principal sensor of DNA
replication stress: ATR is activated at stalled replication forks through
protein-mediated mechanisms (the ATRIP--TOPBP1 axis, and independently
ETAA1) rather than by occupation of its own ATP-binding pocket, and once
active it phosphorylates and activates CHK1, which restrains further origin
firing and protects replication forks from nucleolytic collapse
\citep{kumagai2006,bass2016,saldivar2017}. Genetic and pharmacological
attenuation of ATR or CHK1 activity increases instability at common fragile
sites -- chromosomal loci prone to breakage under replication stress --
consistent with this axis being protective \citep{casper2002,debatisse2012}.

Essentially all existing chemical matter directed at this pathway --
berzosertib, ceralasertib, camonsertib, prexasertib, and related compounds
-- is ATP-competitive kinase inhibitors developed for oncology
synthetic-lethality strategies \citep{fokas2012}. Two cryo-EM structures of
the full human ATR--ATRIP complex released in 2025, bound respectively to
berzosertib (PDB 9L40) and camonsertib (PDB 9L4B) at $\sim$3~{\AA}
resolution, are a substantial improvement over the original 2018 structure
of the complex (PDB 5YZ0, 4.7~{\AA}) and newly resolved that one ATR--ATRIP
complex engages \emph{four} berzosertib molecules -- two in the canonical
ATP-competitive active sites, and two at a previously uncharacterized
allosteric pocket formed at the ATR--ATR dimer interface
\citep{wang2025}. This is a structurally more complex picture than a single
ATP-pocket story, and it raises a practical question for anyone using these
new structures for docking: does a newly deposited, high-resolution
structure automatically make a trustworthy receptor for conventional
ATP-pocket virtual screening?

This is an empirical question, not one that follows from resolution or
recency, and it is the question this pilot study is organized around. Before
screening any candidate against 9L40 or 9L4B, we validate each structure's
suitability as a docking receptor by redocking its own co-crystallized
ligand -- the standard control for docking-protocol accuracy
\citep{trott2010} -- and only proceed to score novel candidates against a
structure that passes. We then report a modest pilot screen of BRICS-derived
\citep{degen2008} candidates against whichever structure(s) pass, as a
demonstration of the validate-before-screening principle applied to a
concrete, currently topical target.

\subsection*{Scope}

This is explicitly a pilot study, not a comprehensive virtual screening
campaign. Docking, by itself, estimates predicted binding and cannot
establish whether a ligand activates, inhibits, or otherwise modulates
ATR or CHK1 -- physiological ATR activation depends on upstream protein
interactions absent from a static structural model
\citep{kumagai2006,bass2016}. All results below are reported as predicted
binding and prioritization, not as demonstrated pharmacology. Molecular
dynamics and MM-GBSA relative energetic estimation, and a broader
cross-kinase selectivity panel, were part of the original design of the
computational pipeline this pilot draws on, but are not included here:
they require substantially more compute (GPU-accelerated MD) and additional
validated receptor structures than were available for this pilot, and are
explicitly left as future work (see Limitations) rather than run at a scale
too small to be informative.

\section*{Methods}

\subsection*{Receptor structures}

Three receptor structures were obtained, in mmCIF format, via the
Registry of Open Data on AWS mirror of the RCSB PDB archive
(\url{s3://pdbsnapshots}, snapshot 2026-01-01), since the primary RCSB PDB
web/FTP endpoints were not reachable from the compute environment used for
this study: PDB 9L40 (ATR kinase domain bound to VE-822/berzosertib,
$\sim$2.9~{\AA}), PDB 9L4B (ATR kinase domain bound to
camonsertib/RP-3500, $\sim$3~{\AA}), and PDB 2YM8 (human CHK1 kinase domain
bound to a chloroisoquinolinylamino-pyrazine-carbonitrile inhibitor,
2.07~{\AA}) \citep{wang2025}. The original 2018 ATR--ATRIP structure
(PDB 5YZ0, 4.7~{\AA} overall) was not used as a docking receptor: at that
resolution, ATP-pocket side-chain positions are not reliably resolved.

For each structure, protein chains were extracted (standard residues only,
via Biopython \citep{cock2009}), and the co-crystallized ligand of the
biologically relevant ATP-competitive active site was identified and
isolated. For 9L40, four copies of the ligand (chemical component
\texttt{A1EIE}) are present per asymmetric unit; contact analysis (heavy
atoms within 4.5~{\AA} of each protein chain) confirmed that residues
A2701/B2701 each contact only their own chain (canonical, single-monomer
ATP-competitive site) while A2702/B2702 contact both chains (the
dimer-interface allosteric site reported by \citet{wang2025}); A2701 was
used as the active-site reference ligand. The analogous single-copy
active-site ligands (chemical component \texttt{A1EIK} for 9L4B;
\texttt{YM8} for 2YM8) were extracted directly. New-style five-character PDB
chemical-component codes overflow the fixed-width residue-name column of
the legacy PDB format and silently corrupt naive PDB-to-PDBQT conversion
(each atom is parsed as an isolated one-atom fragment); ligand records were
therefore rewritten with a placeholder three-character residue name before
conversion. Proteins and ligands were converted to PDBQT with Open Babel
\citep{oboyle2011}, and docking search boxes were centered on each
active-site ligand's heavy-atom centroid (22~{\AA} cubic box).

\subsection*{Seed compound library}

Five compounds formed the seed library (Table~\ref{tab:seeds}): the three
active-site ligands extracted directly from the co-crystal structures above
(berzosertib, camonsertib, and the 2YM8 CHK1 inhibitor), plus prexasertib
and ceralasertib, whose canonical SMILES were obtained via literature search
and independently cross-checked against reported molecular
formula/exact mass (both matched exactly: prexasertib
C$_{18}$H$_{19}$N$_7$O$_2$, exact mass 365.16; ceralasertib
C$_{20}$H$_{24}$N$_6$O$_2$S). The PubChem and ChEMBL REST APIs used by this
pipeline's programmatic seed-retrieval script were not reachable from the
compute environment used for this study, so this smaller, directly-verified
seed set was used in their place; all five structures were confirmed valid
and standardized with RDKit \citep{rdkit} (salt stripping, charge
neutralization, canonicalization).

\begin{table}[htbp]
\centering
\caption{Seed compound library. Compounds marked * were extracted directly from the co-crystallized active-site ligand of the named PDB structure; the remaining SMILES were obtained via literature search and independently verified against reported molecular formula and exact mass.}
\label{tab:seeds}
\small
\begin{tabular}{lll}
\toprule
Compound & Target & Provenance \\
\midrule
Berzosertib (VE-822)   & ATR  & PDB 9L40, active-site ligand* \\
Camonsertib (RP-3500)  & ATR  & PDB 9L4B, active-site ligand* \\
CHK1 inhibitor (YM8)   & CHK1 & PDB 2YM8, active-site ligand* \\
Prexasertib            & CHK1 & Literature SMILES, formula-verified \\
Ceralasertib (AZD6738) & ATR  & Literature SMILES, formula-verified \\
\bottomrule
\end{tabular}
\end{table}

\subsection*{Derivative generation and developability filtering}

The seed set was fragmented with the BRICS algorithm \citep{degen2008} into
22 unique fragments, recombined (\texttt{BRICS.BRICSBuild}, RDKit), and
capped at 5,000 generated candidates (a combinatorial cap, not a natural
stopping point -- see Limitations). Candidates were pre-filtered by a
relaxed Lipinski rule ($\leq$1 violation of molecular weight $\leq$500,
cLogP $\leq$5, H-bond donors $\leq$5, H-bond acceptors $\leq$10)
\citep{lipinski1997}, Veber's rule (rotatable bonds $\leq$10, TPSA
$\leq$140~{\AA}$^2$) \citep{veber2002}, and removal of any structure
matching the RDKit PAINS \citep{baell2010} or Brenk structural-alert
catalogs, leaving 786 candidates. Every retained candidate (plus the 5
seeds) was then scored against four descriptor-based drug-likeness rules
(Lipinski, Veber, Ghose \citep{ghose1999}, Egan \citep{egan2000}); 680 of
791 compounds satisfied $\geq$3/4 rules with no structural alert
(``developability-passed''). All values from this stage are rule-based
predictions, not measurements.

\subsection*{Redocking validation}

Before any candidate was scored, each receptor was validated by extracting
its own active-site ligand, redocking it into the identical receptor/box
with AutoDock Vina \citep{trott2010} (exhaustiveness 32, three independent
seeded replicates: 1000, 1001, 1002), and computing the root-mean-square
deviation (RMSD, heavy atoms, index-matched -- Vina's PDBQT output
preserves input atom order) between the best-scoring pose and the native
crystallographic pose, against a pre-specified pass threshold of
$\leq$2.0~{\AA}, fixed before any receptor was tested. A receptor failing
this control in any replicate is excluded from candidate scoring, with the
reason stated, rather than retried until it passes. Search reproducibility
(pairwise RMSD between replicate best poses, independent of whether that
pose is correct) was reported separately from redocking accuracy (RMSD to
the native pose), since these answer different questions.

\subsection*{Pilot library docking}

The 5 seed compounds plus the top 35 developability-passed derivatives
(ranked by number of rules satisfied, then by molecular weight ascending)
were docked against every receptor structure that passed redocking
validation, using the identical protocol (AutoDock Vina, exhaustiveness 16,
three independent seeded replicates: 2000, 2001, 2002). Ligands were
3D-embedded and protonated at pH 7.4 with Open Babel. The best (most
negative) affinity across replicates is reported per compound, alongside
the mean and standard deviation across replicates.

\subsection*{Computational environment}

All calculations were run with AutoDock Vina 1.2.7 (Python bindings) and
Open Babel (both via pip, CPU only, no GPU available in this environment --
see Limitations regarding molecular dynamics), RDKit 2026.03.6,
and Biopython 1.88. Receptor structures were retrieved
from the RCSB PDB AWS Open Data mirror snapshot dated 2026-01-01. All code
and intermediate data are available at the repository accompanying this
submission (see Availability of data and materials).

\section*{Results}

\subsection*{Redocking validation: one of two ATR structures is unreliable
for docking}

Table~\ref{tab:redocking} summarizes the redocking validation across all
three receptor structures over three independent seeded replicates. 9L4B
(ATR, camonsertib-bound) and 2YM8 (CHK1) both passed with a substantial
margin (mean RMSD $0.61\pm0.02$~{\AA} and $0.90\pm0.01$~{\AA}
respectively, well under the 2.0~{\AA} threshold in every replicate). 9L40
(ATR, berzosertib-bound) failed reproducibly: mean RMSD
$2.96\pm0.06$~{\AA}, exceeding the threshold in all three replicates. This
was not simply search noise -- the pairwise RMSD between 9L40's own
replicate best poses (search reproducibility) was tight ($0.56$~{\AA}
mean), meaning Vina consistently converged on the same, wrong, pose under
this box/protocol, rather than failing to converge at all. Per the
pre-specified rule (Methods), 9L40 was excluded from candidate scoring; the
pilot library docking below uses 9L4B only.

\begin{table}[htbp]
\centering
\caption{Redocking validation across three independent seeded Vina replicates (exhaustiveness 32).}
\label{tab:redocking}
\begin{tabular}{lccc}
\toprule
Structure & Mean RMSD ({\AA}) & Passes $\leq$2.0~{\AA} & Search reproducibility ({\AA}) \\
\midrule
9L40 (ATR, VE-822)        & $2.96 \pm 0.06$ & \textbf{No} (all 3 replicates) & $0.56$ \\
9L4B (ATR, camonsertib)   & $0.61 \pm 0.02$ & Yes (all 3 replicates)         & $0.045$ \\
2YM8 (CHK1)               & $0.90 \pm 0.01$ & Yes (all 3 replicates)         & $0.078$ \\
\bottomrule
\end{tabular}
\end{table}

\subsection*{Seed and derivative library}

The 5-compound seed set yielded 22 unique BRICS fragments, from which 5,000
recombination candidates were generated (capped) and 786 passed the initial
Lipinski/Veber/PAINS/Brenk pre-filter. Of the resulting 791 compounds (786
derivatives + 5 seeds), 680 (86\%) satisfied $\geq$3/4 developability rules
with no PAINS/Brenk structural alert.

\subsection*{Pilot docking against the validated ATR structure}

The 40-compound pilot set (5 seeds + 35 developability-passed derivatives)
was docked against 9L4B (exhaustiveness 16, three seeded replicates each).
Best-pose predicted affinity across the set ranged from $-9.80$ to
$-7.82$~kcal/mol (mean $-8.66\pm0.45$~kcal/mol across all 40 compounds).
Table~\ref{tab:docking} shows the top 10 ranked compounds. Camonsertib --
the receptor's own co-crystallized ligand -- ranked first, as expected for
a structure the receptor was solved with. No BRICS derivative exceeded the
best seed compound's predicted affinity; the closest, \texttt{deriv\_002}
($-9.47$~kcal/mol), came within $0.33$~kcal/mol of camonsertib. The three
top-ranked derivatives (\texttt{deriv\_002}, \texttt{deriv\_003},
\texttt{deriv\_019}) share camonsertib's oxabicyclic-ether/pyrazolopyrazine
core with a single substituent varied (methylamino, methoxy, and
dimethylamino respectively, in place of camonsertib's own morpholine
substituent) -- an interpretable and expected outcome of BRICS
recombination re-using the best-scoring seed compound's own fragments in a
new context, rather than an independent discovery of novel chemical matter.

\begin{table}[htbp]
\centering
\caption{Top 10 ranked compounds from the 40-compound pilot docking against 9L4B (mean $\pm$ SD and best across 3 seeded replicates, exhaustiveness 16).}
\label{tab:docking}
\small
\begin{tabular}{llcc}
\toprule
Compound & Source & Mean affinity (kcal/mol) & Best affinity (kcal/mol) \\
\midrule
camonsertib (seed) & seed       & $-9.78 \pm 0.01$ & $-9.80$ \\
deriv\_002          & derivative & $-9.46 \pm 0.01$ & $-9.47$ \\
ceralasertib (seed) & seed       & $-9.39 \pm 0.01$ & $-9.40$ \\
deriv\_003          & derivative & $-9.33 \pm 0.01$ & $-9.34$ \\
deriv\_019          & derivative & $-9.25 \pm 0.03$ & $-9.29$ \\
deriv\_032          & derivative & $-9.06 \pm 0.17$ & $-9.15$ \\
deriv\_022          & derivative & $-9.08 \pm 0.02$ & $-9.10$ \\
berzosertib (seed)  & seed       & $-8.94 \pm 0.05$ & $-8.99$ \\
deriv\_033          & derivative & $-8.96 \pm 0.02$ & $-8.98$ \\
deriv\_026          & derivative & $-8.94 \pm 0.00$ & $-8.94$ \\
\bottomrule
\end{tabular}
\end{table}

\section*{Discussion}

The central finding of this pilot is methodological rather than chemical:
\textbf{a newly deposited, higher-resolution structure is not automatically
a reliable docking receptor, and only an explicit redocking control
establishes this, per structure, before any candidate is trusted.} 9L40 and
9L4B were solved in the same study, at comparable resolution
\citep{wang2025}, yet only one reproducibly passed a standard accuracy
control under an otherwise identical protocol. A screening campaign that
skipped this control and used 9L40 by default (e.g., because it was the
first hit in a search, or because it captures the more novel allosteric
site) would have scored every candidate against a receptor conformation the
docking protocol itself cannot reliably reproduce. Whether the 9L40 failure
reflects a genuine limitation of rigid-receptor Vina docking against this
particular conformational state (e.g., an induced-fit effect specific to
the dimer-interface-engaged assembly) or a more mundane box/protocol
mismatch is not resolved by this pilot and is worth follow-up, but the
practical implication for anyone reusing these structures is the same
either way: validate before you screen.

The pilot docking result is, appropriately for a 35-derivative pilot,
modest rather than dramatic: no generated derivative exceeded the predicted
affinity of the seed compounds, and the best-performing derivatives are
recognizable as close analogs of camonsertib rather than structurally novel
chemotypes. This is consistent with a BRICS recombination run over a
5-compound seed set and a fragment pool of only 22 pieces: the search space
explored is inherently local to the seed chemistry, and a pilot at this
scale should not be expected to surface a structurally distinct chemotype
with superior predicted affinity. What this pilot does show is that the
pipeline behaves sensibly at small scale -- it recovers the seed compound's
own high rank (unsurprising, since camonsertib is the receptor's native
ligand, but a useful sanity check that the docking protocol is not
producing arbitrary rankings) and produces derivatives whose structure-score
relationship is interpretable (small, defined substituent changes at one
position, with correspondingly small score changes) rather than
erratic. Whether a larger seed set, a larger fragment pool, or multiple
recombination rounds would surface derivatives that improve on the seed
compounds is an open question this pilot's scale cannot answer.

\subsection*{Limitations}

This is a small pilot, and its scope was set by what could be executed
directly and verifiably, not by what would be most impressive to report.
Specifically: (1) the BRICS recombination cap (5,000) is combinatorial, not
principled, and a larger or differently seeded run would sample a
different, likely larger, region of the accessible chemical space; (2) the
40-compound pilot library is a small, deterministic subset of the
developability-passed derivatives (ranked by rule count then molecular
weight), not a comprehensive screen, and enrichment-style retrieval
statistics (ROC-AUC, enrichment factor) are not meaningful at this scale and
are not reported; (3) no molecular dynamics or MM-GBSA relative energetic
estimation was run -- both require substantially more compute (GPU
acceleration for MD at any informative timescale) than was available, and
docking scores alone are known to correlate imperfectly with binding
affinity; (4) no cross-kinase selectivity panel was run -- this would
require additional receptor structures, each independently redocking-validated
by the same standard applied to the three structures here, which was beyond
this pilot's scope; (5) as throughout, nothing here establishes ATR/CHK1
functional modulation, inhibition, or activation -- only predicted binding
and ranking; (6) the seed and reference SMILES for prexasertib and
ceralasertib were obtained via web search rather than a structured chemical
database query (PubChem/ChEMBL were not reachable) and, while independently
verified against reported molecular formula and exact mass, should be
re-confirmed against a primary chemical database before further use.

\section*{Conclusions}

Redocking validation -- not resolution, recency, or biological interest --
is what determines whether a deposited structure is a trustworthy docking
receptor. Applied to two contemporaneous, comparable-resolution ATR
structures, this control passed one and failed the other, a distinction
that would not have been visible without running it. Screened against the
validated structure, a small BRICS-derived pilot library reproduced the
expected top rank of the receptor's own native ligand and generated
interpretable, if modest, analogs of it, without yet surfacing a chemotype
that improves on the seed compounds -- a realistic, honestly reported
outcome for a pilot of this scale, and one further reason to treat larger
claims from similarly small screens with caution. All structures, code, and
intermediate data used in this study are openly available for independent
verification and reuse.

\section*{Availability of data and materials}
All code, the seed and derivative compound tables, redocking validation
results, and pilot docking results are available at
\url{https://github.com/ethanyyy123/research-molecule-editing}
(branch \url{claude/cadd-fragile-sites-optimization-0u9fsp} at the time
of this submission; an archived, versioned release should be minted, e.g.
via Zenodo, before final submission). Receptor structures are publicly
available from RCSB PDB (9L40, 9L4B, 2YM8) and were retrieved via the AWS
Open Data mirror (\url{s3://pdbsnapshots}).

\section*{Competing interests}
[PLACEHOLDER -- e.g. ``The authors declare no competing interests,'' if accurate; confirm before submission.]

\section*{Funding}
[PLACEHOLDER -- funding source(s), or ``Not applicable'' if self-funded/no external support.]

\section*{Authors' contributions}
[PLACEHOLDER -- e.g. following CRediT roles: conceptualization, methodology, software, formal analysis, investigation, writing.]

\section*{Acknowledgements}
[PLACEHOLDER -- optional; e.g. acknowledgement of compute resources or advisors, if applicable.]

\bibliography{refs}

\end{document}
